% Part: normal-modal-logic
% Chapter: tableaux
% Section: soundness
%READERNOTE{SOL6-A176: निर्दोषतेच्या सिद्धतेत मोडल प्रतिमान आणि मोडल पूर्तीची संज्ञा वापरली. नियमानुसार तयार झालेल्या शाखेची पूर्वचिन्हे आरंभीच्या रिक्त नसलेल्या खंडांखाली बंद असतात. म्हणून नवे पूर्वचिन्ह वाढवताना त्याचे आधीपासूनचे वंशज नसतात आणि फक्त पालकापासून नव्या पूर्वचिन्हापर्यंतची प्राप्यता तपासणे पुरेसे असते.}

\documentclass[../../../include/open-logic-section]{subfiles}

\begin{document}

\olfileid{nml}{tab}{sou}

\olsection{\Log{K} साठी निर्दोषता}

\begin{editorial}
  निर्दोषतेचा हा पुरावा अभिजात विधानीय तर्कशास्त्राचा
  निर्दोषतेचा पुरावा पुन्हा मांडतो; म्हणजे सर्व काही
  सुरुवातीपासून सिद्ध करतो. स्वयंपूर्ण पुरावा हवा
  असेल, तर हे योग्य आहे. नेहमीच्या टॅब्लोची निर्दोषता
  तुम्ही आधी पाहिली असेल, तर येथे पुनरुक्ती वाटेल.
  पुढे स्वयंपूर्ण आवृत्ती आणि अमोडल प्रकरणावर आधारलेली
  आवृत्ती यांपैकी एक निवडता यावी, अशी योजना आहे.
\end{editorial}

\begin{explain}
  पूर्वचिन्हयुक्त !!{tableau}s निर्दोष आहेत हे
  दाखवण्यासाठी, पुढील गृहीतकांचा
  \[
  \sFmla{\True}{!B_1}[1], \dots, \sFmla{\True}{!B_n}[1], \sFmla{\False}{!A}[1]
  \]
  बंद !!{tableau} असेल, तर $!B_1, \dots, !B_n
  \Entails !A$ हे सिद्ध करावे लागेल. व्यत्यास सिद्ध
  करणे सोपे आहे: एखाद्या $\mModel{M}$ मध्ये
  $w$ या जगात सर्व $i=1$, \dots,~$n$ साठी
  $\mSat{M}{!B_i}[w]$, पण $\mSat/{M}{!A}[w]$
  असेल, तर कोणताही !!{tableau} बंद होऊ शकत नाही.
  अशा प्रतिदृष्टांत प्रतिमानावरून !!{tableau} ची सुरुवातीची
  गृहीतके पूर्ततायोग्य आहेत हे दिसते. पुराव्याची युक्ती
  अशी: !!a{tableau} च्या शाखेवरील सर्व
  पूर्वचिन्हयुक्त !!{formula}s पूर्ततायोग्य असतील, तर
  कोणताही नियम लावल्यावर तयार होणाऱ्या विस्तारित
  शाखांपैकी किमान एक शाखा पूर्ततायोग्य असते. बंद शाखा
  अपूर्ततायोग्य असतात. म्हणून पूर्वचिन्हयुक्त !!{formula}s
  च्या पूर्ततायोग्य संचासाठीच्या कोणत्याही !!{tableau}
  मध्ये किमान एक खुली शाखा असली पाहिजे.

  मोडल प्रकरणात ही युक्ती वापरण्यासाठी ‘पूर्ततायोग्य’
  ही व्याख्या मोडल प्रतिमाने आणि पूर्वचिन्हे यांच्यापर्यंत
  विस्तारित करावी लागते. त्यानंतर पुरावा सरळ आहे.
\end{explain}

\begin{defn}
  $P$ हा पूर्वचिन्हांचा एखादा संच असू द्या; म्हणजे
  $P \subseteq (\PosInt)^* \setminus
  \{\emptyseq\}$. तसेच $\mModel{M}$ हे प्रतिमान असू द्या.
  $\sigma$ आणि $\sigma.n$ ही दोन्ही~$P$ मध्ये असताना
  नेहमी $Rf(\sigma)f(\sigma.n)$ असेल, तर
  $f\colon P \to W$ या फलनाला $\mModel{M}$
  मधील~$P$ चे \emph{अर्थनिर्धारण} म्हणतात.

  पूर्वचिन्हांच्या~$P$ या संचाच्या अर्थनिर्धारणाच्या
  संदर्भात पुढील व्याख्या करता येतात:
  \begin{enumerate}
  \item $\mModel{M}$ हे $\sFmla{\True}{!A}[\sigma]$
    पूर्ण करते तेव्हा आणि तेव्हाच
    $\mSat{M}{!A}[f(\sigma)]$.
  \item $\mModel{M}$ हे $\sFmla{\False}{!A}[\sigma]$
    पूर्ण करते तेव्हा आणि तेव्हाच
    $\mSat/{M}{!A}[f(\sigma)]$.
  \end{enumerate}
\end{defn}

\begin{defn}
  $\Gamma$ हा पूर्वचिन्हयुक्त !!{formula}s चा
  संच असू द्या आणि त्यात येणाऱ्या पूर्वचिन्हांचा
  संच $P(\Gamma)$ असू द्या. $f$ हे $\mModel{M}$
  मधील~$P(\Gamma)$ चे अर्थनिर्धारण असेल, तर
  $\mModel{M}$ ने~$f$ च्या संदर्भात~$\Gamma$
  पूर्ण करणे, म्हणजे $\mSat{M}{\Gamma}[f]$,
  याचा अर्थ $\mModel{M}$ ने~$f$ च्या संदर्भात
  $\Gamma$ मधील प्रत्येक पूर्वचिन्हयुक्त !!{formula}
  पूर्ण करणे असा आहे. $\mSat{M}{\Gamma}[f]$
  होईल असे एखादे प्रतिमान~$\mModel{M}$ आणि
  $P(\Gamma)$ चे अर्थनिर्धारण~$f$ अस्तित्वात असेल,
  तर आणि तरच $\Gamma$ हा संच \emph{पूर्ततायोग्य} असतो.
\end{defn}

\begin{prop}
  एखाद्या !!{formula}~$!A$ आणि पूर्वचिन्ह~$\sigma$
  यांसाठी $\Gamma$ मध्ये
  $\sFmla{\True}{!A}[\sigma]$ आणि
  $\sFmla{\False}{!A}[\sigma]$ ही दोन्ही असतील,
  तर $\Gamma$ अपूर्ततायोग्य आहे.
\end{prop}

\begin{proof}
  $\mSat{M}{!A}[f(\sigma)]$ आणि
  $\mSat/{M}{!A}[f(\sigma)]$ ही दोन्ही सत्य होतील
  असे प्रतिमान~$\mModel{M}$ आणि $P(\Gamma)$ चे
  अर्थनिर्धारण~$f$ असू शकत नाही.
\end{proof}

\begin{thm}[निर्दोषता]
  \ollabel{thm:tableau-soundness}
  $\Gamma$ साठी बंद !!{tableau} असेल, तर
  $\Gamma$ अपूर्ततायोग्य आहे.
\end{thm}

\begin{proof}
!!a{tableau} वरील एखादी शाखा पूर्ततायोग्य असणे म्हणजे
तिच्यावरील !!{signed formula}s चा संच पूर्ततायोग्य असणे.
आणि !!a{tableau} मध्ये किमान एक पूर्ततायोग्य शाखा
असेल, तर त्याला पूर्ततायोग्य म्हणू.
येथे नियमानुसार तयार केलेल्या टॅब्लोंचाच विचार करतो. सुरुवातीच्या
गृहीतकांचे पूर्वचिन्ह $1$ असते आणि प्रत्येक नवे पूर्वचिन्ह आधी
वापरलेल्या पूर्वचिन्हाच्या शेवटी एक संख्या जोडून तयार होते.
म्हणून शाखेवरील पूर्वचिन्हांचा संच प्रत्येक रिक्त नसलेल्या आरंभीच्या
खंडाखाली बंद असतो.

आपण पुढील विधान दाखवू: पूर्ततायोग्य !!{tableau} ला
अनुमानाचा कोणताही नियम लावून विस्तारित केले, तरी
मिळणारा !!{tableau} पूर्ततायोग्य असतो. यावरून प्रमेय
सिद्ध होईल. कारण बंद !!{tableau} हा
सुरुवातीला फक्त~$\Gamma$ मधील गृहीतके असलेल्या
!!{tableau} ला अनुमानाचे नियम लावत तयार होतो.
म्हणून $\Gamma$ पूर्ततायोग्य असेल, तर त्यासाठीचा
प्रत्येक !!{tableau} पूर्ततायोग्य राहील. पण बंद
!!{tableau} पूर्ततायोग्य असू शकत नाही: त्याच्या सर्व
शाखा बंद असतात, आणि बंद शाखा अपूर्ततायोग्य असतात.

आपल्याकडे पूर्ततायोग्य !!{tableau} आहे असे समजू;
म्हणजे किमान एक पूर्ततायोग्य शाखा असलेला !!a{tableau}
आहे. अनुमानाचा
नियम लावल्याने शाखेवर !!{signed formula}s
जोडली जातात किंवा शाखेच्या दोन शाखा फुटतात.
नियमामुळे न बदलणारी एखादी पूर्ततायोग्य शाखा
!!{tableau} मध्ये असेल, तर विस्तारित
!!{tableau} मध्येही ती पूर्ततायोग्य राहते. म्हणून
नियम पूर्ततायोग्य शाखेला लावला जातो त्या
प्रकरणाचा विचार पुरेसा आहे.

त्या शाखेवरील !!{signed formula}s चा संच
$\Gamma$ असू द्या आणि नियम
$\sFmla{S}{!A}[\sigma] \in \Gamma$ या
!!{signed formula} ला लावला जातो असे समजू.
नियमाने शाखा फुटत नसेल, तर $\Gamma$ मध्ये
नियमाचे निष्कर्ष जोडून मिळालेली विस्तारित
शाखा अजूनही पूर्ततायोग्य आहे हे दाखवावे लागते.
शाखा फुटत असेल, तर मिळणाऱ्या दोन शाखांपैकी
किमान एक पूर्ततायोग्य आहे हे दाखवावे लागते.
\tagfalse{prvDiamond}
प्रथम शाखा न फुटणाऱ्या नियमांचा विचार करू.
\begin{enumerate}
\item $\sFmla{\True}{\lnot !B}[\sigma] \in \Gamma$
  याला
  $\TRule{\True}{\lnot}$ लावून शाखा विस्तारित
  केली. आता शाखेवरील !!{signed formula}s चा
  संच $\Gamma \cup
  \{\sFmla{\False}{!B}[\sigma]\}$ आहे.
  $\mSat{M}{\Gamma}[f]$ असे समजू. विशेषतः
  $\mSat{M}{\lnot !B}[f(\sigma)]$; म्हणून
  $\mSat/{M}{!B}[f(\sigma)]$. म्हणजे
  $\mModel{M}$ हे~$f$ च्या संदर्भात
  $\sFmla{\False}{!B}[\sigma]$ पूर्ण करते.
\item $\sFmla{\False}{\lnot !B}[\sigma] \in \Gamma$
  याला
  $\TRule{\False}{\lnot}$ लावून शाखा विस्तारित
  केली: सरावासाठी.
\item $\sFmla{\True}{!B \land !C}[\sigma] \in \Gamma$
  याला
  $\TRule{\True}{\land}$ लावून शाखा विस्तारित
  केली. यात शाखेवर
  $\sFmla{\True}{!B}[\sigma]$ आणि
  $\sFmla{\True}{!C}[\sigma]$ ही दोन नवी
  !!{signed formula}s मिळतात.
  $\mSat{M}{\Gamma}[f]$ असे समजू; विशेषतः
  $\mSat{M}{!B \land !C}[f(\sigma)]$.
  म्हणून $\mSat{M}{!B}[f(\sigma)]$ आणि
  $\mSat{M}{!C}[f(\sigma)]$. याचा अर्थ
  $\mModel{M}$ हे~$f$ च्या संदर्भात
  $\sFmla{\True}{!B}[\sigma]$ आणि
  $\sFmla{\True}{!C}[\sigma]$ ही दोन्ही
  पूर्ण करते.
\item $\sFmla{\False}{!B \lor !C}[\sigma] \in \Gamma$
  याला
  $\TRule{\False}{\lor}$ लावून शाखा विस्तारित
  केली: सरावासाठी.
\item $\sFmla{\False}{!B \lif !C}[\sigma] \in \Gamma$
  याला
  $\TRule{\False}{\lif}$ लावून शाखा विस्तारित
  केली. यात शाखेवर
  $\sFmla{\True}{!B}[\sigma]$ आणि
  $\sFmla{\False}{!C}[\sigma]$ ही दोन नवी
  !!{signed formula}s मिळतात.
  $\mSat{M}{\Gamma}[f]$ असे समजू; विशेषतः
  $\mSat/{M}{!B \lif !C}[f(\sigma)]$.
  म्हणून $\mSat{M}{!B}[f(\sigma)]$ आणि
  $\mSat/{M}{!C}[f(\sigma)]$. याचा अर्थ
  $\mModel{M}, f$ हे
  $\sFmla{\True}{!B}[\sigma]$ आणि
  $\sFmla{\False}{!C}[\sigma]$ ही दोन्ही
  पूर्ण करते.

\iftag{prvBox}{%
\item $\sFmla{\True}{\Box !B}[\sigma] \in \Gamma$
  याला
  $\TRule{\True}{\Box}$ लावून शाखा विस्तारित केली:
  \iftag{probBox}{सरावासाठी.}{शाखेवर
    $\sFmla{\True}{!B}[\sigma.n]$ हे नवे
    !!{signed
      formula} मिळते. येथे
    $\sigma.n \in P(\Gamma)$, कारण
    $\sigma.n$ वापरलेले असले पाहिजे.
    $\mSat{M}{\Gamma}[f]$ असे समजू; विशेषतः
    $\mSat{M}{\Box !B}[f(\sigma)]$.
    $f$ हे पूर्वचिन्हांचे अर्थनिर्धारण आहे आणि
    $\sigma$, $\sigma.n \in P(\Gamma)$,
    म्हणून $Rf(\sigma)f(\sigma.n)$.
    त्यामुळे $\mSat{M}{!B}[f(\sigma.n)]$;
    म्हणजे $\mModel{M}, f$ हे
    $\sFmla{\True}{!B}[\sigma.n]$ पूर्ण करते.}

\item $\sFmla{\False}{\Box !B}[\sigma] \in \Gamma$
  याला
  $\TRule{\False}{\Box}$ लावून शाखा विस्तारित केली:
  \iftag{probBox}{सरावासाठी.}{शाखेवर
    $\sFmla{\False}{!B}[\sigma.n]$ हे नवे
    !!{signed
      formula} मिळते. येथे
    $\sigma.n$ हे शाखेवरील नवे पूर्वचिन्ह,
    म्हणजे $\sigma.n \notin P(\Gamma)$.
    $\Gamma$ पूर्ततायोग्य असल्याने,
    $\mSat{M}{\Gamma}[f]$ होईल असे
    $\mModel{M}$ आणि $P(\Gamma)$ चे
    अर्थनिर्धारण~$f$ असते; विशेषतः
    $\mSat/{M}{\Box !B}[f(\sigma)]$.
    $\Gamma \cup
    \{\sFmla{\False}{!B}[\sigma.n]\}$
    पूर्ततायोग्य आहे हे दाखवायचे आहे.
    यासाठी $P(\Gamma) \cup \{\sigma.n\}$
    चे अर्थनिर्धारण पुढीलप्रमाणे ठरवू:

  $\mSat/{M}{\Box !B}[f(\sigma)]$ असल्याने,
  $Rf(\sigma)w$ आणि $\mSat/{M}{!B}[w]$
  होईल असे $w \in W$ असते. $f'$ हे~$f$
  प्रमाणेच असू द्या, फक्त
  $f'(\sigma.n) = w$ असे ठरवा.
  $f'(\sigma) = f(\sigma)$ आणि
  $Rf(\sigma)w$ असल्याने
  $Rf'(\sigma)f'(\sigma.n)$.
  पूर्वचिन्हांच्या आरंभीच्या खंडांखालील बंदतेमुळे नव्या
  $\sigma.n$ पासून वाढवलेले कोणतेही पूर्वचिन्ह आधी वापरलेले
  नसते. त्यामुळे अर्थनिर्धारणासाठी तपासायची नवी जोडी फक्त
  $(\sigma,\sigma.n)$ हीच आहे; म्हणून $f'$
  हे $P(\Gamma) \cup \{\sigma.n\}$ चे
  अर्थनिर्धारण आहे. स्पष्टपणे
  $\mSat/{M}{!B}[f'(\sigma.n)]$.
  प्रत्येक $\sigma' \in P(\Gamma)$ साठी
  $f(\sigma') = f'(\sigma')$ असल्याने
  $\mSat{M}{\Gamma}[f']$. म्हणून
  $\mModel{M}, f'$ हे $\Gamma \cup
  \{\sFmla{\False}{!B}[\sigma.n]\}$ पूर्ण करते.}
  }{}

\iftag{prvDiamond}{%
\item $\sFmla{\True}{\Diamond !B}[\sigma] \in \Gamma$
  याला
  $\TRule{\True}{\Diamond}$ लावून शाखा विस्तारित
  केली: \iftag{probDiamond}{सरावासाठी.}{शाखेवर
    $\sFmla{\True}{!B}[\sigma.n]$ हे नवे
    !!{signed
      formula} मिळते. येथे
    $\sigma.n$ हे शाखेवरील नवे पूर्वचिन्ह,
    म्हणजे $\sigma.n \notin P(\Gamma)$.
    $\Gamma$ पूर्ततायोग्य असल्याने,
    $\mSat{M}{\Gamma}[f]$ होईल असे
    $\mModel{M}$ आणि $P(\Gamma)$ चे
    अर्थनिर्धारण~$f$ असते; विशेषतः
    $\mSat{M}{\Diamond !B}[f(\sigma)]$.
    $\Gamma \cup
    \{\sFmla{\True}{!B}[\sigma.n]\}$
    पूर्ततायोग्य आहे हे दाखवायचे आहे.
    यासाठी $P(\Gamma) \cup \{\sigma.n\}$
    चे अर्थनिर्धारण पुढीलप्रमाणे ठरवू:

  $\mSat{M}{\Diamond !B}[f(\sigma)]$ असल्याने,
  $Rf(\sigma)w$ आणि $\mSat{M}{!B}[w]$
  होईल असे $w \in W$ असते. $f'$ हे~$f$
  प्रमाणेच असू द्या, फक्त
  $f'(\sigma.n) = w$ असे ठरवा.
  $f'(\sigma) = f(\sigma)$ आणि
  $Rf(\sigma)w$ असल्याने
  $Rf'(\sigma)f'(\sigma.n)$.
  पूर्वचिन्हांच्या आरंभीच्या खंडांखालील बंदतेमुळे नव्या
  $\sigma.n$ पासून वाढवलेले कोणतेही पूर्वचिन्ह आधी वापरलेले
  नसते. त्यामुळे अर्थनिर्धारणासाठी तपासायची नवी जोडी फक्त
  $(\sigma,\sigma.n)$ हीच आहे; म्हणून $f'$
  हे $P(\Gamma) \cup \{\sigma.n\}$ चे
  अर्थनिर्धारण आहे. स्पष्टपणे
  $\mSat{M}{!B}[f'(\sigma.n)]$.
  प्रत्येक $\sigma' \in P(\Gamma)$ साठी
  $f(\sigma') = f'(\sigma')$ असल्याने
  $\mSat{M}{\Gamma}[f']$. म्हणून
  $\mModel{M}, f'$ हे $\Gamma \cup
  \{\sFmla{\True}{!B}[\sigma.n]\}$ पूर्ण करते.}
\item $\sFmla{\False}{\Diamond !B}[\sigma] \in \Gamma$
  याला
  $\TRule{\False}{\Diamond}$ लावून शाखा विस्तारित
  केली: \iftag{probDiamond}{सरावासाठी.}{शाखेवर
    $\sFmla{\False}{!B}[\sigma.n]$ हे नवे
    !!{signed
      formula} मिळते. येथे
    $\sigma.n \in P(\Gamma)$, कारण
    $\sigma.n$ वापरलेले असले पाहिजे.
    $\mSat{M}{\Gamma}[f]$ असे समजू; विशेषतः
    $\mSat/{M}{\Diamond !B}[f(\sigma)]$.
    $f$ हे पूर्वचिन्हांचे अर्थनिर्धारण आहे आणि
    $\sigma$, $\sigma.n \in P(\Gamma)$,
    म्हणून $Rf(\sigma)f(\sigma.n)$.
    त्यामुळे $\mSat/{M}{!B}[f(\sigma.n)]$;
    म्हणजे $\mModel{M}, f$ हे
    $\sFmla{\False}{!B}[\sigma.n]$ पूर्ण करते.}
  }{}
\end{enumerate}
आता शाखा फुटणाऱ्या नियमांचा विचार करू.
\begin{enumerate}
\item $\sFmla{\False}{!B \land !C}[\sigma] \in \Gamma$
  याला
  $\TRule{\False}{\land}$ लावून शाखा विस्तारित
  केली. यात दोन शाखा मिळतात: डावीकडची
  $\sFmla{\False}{!B}[\sigma]$ मधून, आणि
  उजवीकडची $\sFmla{\False}{!C}[\sigma]$
  मधून पुढे जाते. $\mSat{M}{\Gamma}[f]$
  असे समजू; विशेषतः
  $\mSat/{M}{!B \land !C}[f(\sigma)]$.
  मग $\mSat/{M}{!B}[f(\sigma)]$ किंवा
  $\mSat/{M}{!C}[f(\sigma)]$.
  पहिल्या प्रकरणात $\mModel{M}, f$ हे
  $\sFmla{\False}{!B}[\sigma]$ पूर्ण करते,
  म्हणून डावीकडची शाखा पूर्ततायोग्य असते.
  दुसऱ्या प्रकरणात $\mModel{M}, f$ हे
  $\sFmla{\False}{!C}[\sigma]$ पूर्ण करते,
  म्हणून उजवीकडची शाखा पूर्ततायोग्य असते.
\item $\sFmla{\True}{!B \lor !C}[\sigma] \in \Gamma$
  याला
  $\TRule{\True}{\lor}$ लावून शाखा विस्तारित
  केली: सरावासाठी.
\item $\sFmla{\True}{!B \lif !C}[\sigma] \in \Gamma$
  याला
  $\TRule{\True}{\lif}$ लावून शाखा विस्तारित
  केली: सरावासाठी.
\end{enumerate}
\end{proof}

\begin{prob}
\olref[nml][tab][sou]{thm:tableau-soundness} चा
पुरावा पूर्ण करा.
\end{prob}

\begin{cor}
\ollabel{cor:entailment-soundness}
$\Gamma \Proves !A$ असेल, तर
$\Gamma \Entails !A$.
\end{cor}

\begin{proof}
  $\Gamma \Proves !A$ असेल, तर
  $!B_1$, \dots, $!B_n \in \Gamma$ अशी काही सूत्रे
  असतात की $\Delta =
  \{\sFmla{\False}{!A}[1], \sFmla{\True}{!B_1}[1],
  \dots, \sFmla{\True}{!B_n}[1]\}$ साठी
  बंद !!{tableau} असतो.
  $\Gamma \Entails !A$ हे दाखवायचे आहे.
  तसे नाही असे समजू. मग एखाद्या
  $\mModel{M}$ आणि~$w$ साठी सर्व
  $i=1$, \dots,~$n$ बाबत
  $\mSat{M}{!B_i}[w]$, पण
  $\mSat/{M}{!A}[w]$. $f(1) = w$
  असे ठेवा. मग $f$ हे~$\mModel{M}$ मधील
  $P(\Delta)$ चे अर्थनिर्धारण आहे आणि
  $\mModel{M}$ हे~$f$ च्या संदर्भात
  $\Delta$ पूर्ण करते. पण
  \olref{thm:tableau-soundness} नुसार
  $\Delta$ साठी बंद !!{tableau} असल्यामुळे
  तो अपूर्ततायोग्य आहे. हा विरोधाभास आहे.
  म्हणून $\Gamma \Entails !A$.
\end{proof}

\begin{cor}
\ollabel{cor:weak-soundness}
$\Proves !A$ असेल, तर $!A$ सर्व
प्रतिमानांमध्ये सत्य आहे.
\end{cor}

\end{document}
